%% Finitary Boolean operators, canonical extensions, and descriptive duality.

\subsection{Boolean operators and their relational representation}
\label{sec:p10-boolean-operators}

This development uses the held 19 October 2005 preprint of Venema's
\emph{Algebras and Coalgebras}; its local page and result numbers
are cited as preprint locators \cite{Venema2005Preprint}.
The published 2007 chapter has its own bibliography entry
\cite{Venema2007}.
For the published literature on varieties of complex algebras,
see also \textcite{Goldblatt1989}.
The general representation theorem is attributed to
\textcite{JonssonTarski1951}.
The proofs below assume \(\mathsf{ZF}\) plus BPIT, as in the
preceding Stone construction.

\begin{hbdefinition}{def:p10-bao}{Boolean algebra with operators}
Fix a set \(\Lambda\) of additional operation symbols, with a finite
arity \(n_f\in\mathbb N\) for every \(f\in\Lambda\).
A Boolean algebra with these operators is
\(\mathbf B=(B,(f^B)_{f\in\Lambda})\), whose Boolean reduct is \(B\),
such that for every positive-arity \(f\) and each coordinate,
\[
\begin{aligned}
 f^B(b_1,\ldots,0,\ldots,b_n)&=0,\\
 f^B(b_1,\ldots,a\vee c,\ldots,b_n)
 &=f^B(b_1,\ldots,a,\ldots,b_n)
       \vee f^B(b_1,\ldots,c,\ldots,b_n).
\end{aligned}
\]
These are coordinatewise normality and additivity.
They imply monotonicity in each coordinate.
For \(n_f=0\), \(f^B()\) is a designated Boolean element; the
coordinate conditions are vacuous and do not require it to be zero.
A homomorphism preserves the Boolean operations and every \(f^B\),
including nullary ones.

A normal modal algebra has one unary operator \(\Diamond\).
Its dual is \(\Box a=\neg\Diamond\neg a\), which preserves finite
meets and \(1\).
The general definition allows an arbitrary set of symbols of
different finite arities
\cite[preprint Definition~3.11, p.~9]{Venema2005Preprint}.
\index{Boolean algebra with operators}\index{modal algebra}
\end{hbdefinition}

\begin{hbdefinition}{def:p10-bao-complex-algebra}{Relational frames and complex algebras}
A relational \(\Lambda\)-frame
\(\mathfrak F=(W,(R_f)_{f\in\Lambda})\) has
\(R_f\subseteq W^{n_f+1}\).
Empty carriers are permitted in this algebraic construction, so
the trivial Boolean algebra can be represented. A pointed modal
model still requires a chosen point.
The complex algebra \(\mathfrak F^+\) has Boolean reduct
\(\mathcal P(W)\) and operations
\[
 f^{\mathfrak F^+}(A_1,\ldots,A_n)
 =\{w:\exists w_1\in A_1\cdots\exists w_n\in A_n\
                     R_f(w,w_1,\ldots,w_n)\}.
\]
For \(n=0\), this is \(\{w:R_f(w)\}\).
For positive arity these operations preserve arbitrary unions
in every coordinate, including empty unions.
This is the complex-operator construction of
\cite[preprint Definition~3.5, p.~8]{Venema2005Preprint},
with the empty-frame convention explicitly included here.
\end{hbdefinition}

\begin{hbdefinition}{def:p10-bao-ultrafilter-frame}{The ultrafilter frame of a Boolean algebra with operators}
Let \(X_B\) be the Stone space of the Boolean reduct.
Define \(\mathbf B^\bullet=(X_B,(R_f)_{f\in\Lambda})\) by
\[
 R_f(U,V_1,\ldots,V_n)
 \quad\Longleftrightarrow\quad
 \forall b_1\in V_1\cdots\forall b_n\in V_n\
                      f^B(b_1,\ldots,b_n)\in U.
\]
For a nullary symbol, this reads \(R_f(U)\) iff \(f^B()\in U\).
No witness choices enter the definition.
The source calls this the ultrafilter frame or canonical structure
\cite[preprint Definition~5.12, p.~19]{Venema2005Preprint}.
\end{hbdefinition}

\begin{hbexample}{def:p10-bao-binary-nullary}{Binary meet and a nullary operator}
For the binary operator \(f(a,b)=a\wedge b\), the canonical relation is
\[
 R_f(U,V_1,V_2)\quad\Longleftrightarrow\quad U=V_1=V_2.
\]
Indeed, setting the other argument to \(1\) in the defining relation
gives \(V_i\subseteq U\), hence equality of the ultrafilters.
Conversely, a single ultrafilter is closed under meets.
The complex operation on subsets of \(X_B\) is therefore intersection.

For a nullary operator with value \(c\in B\), the relation is
\(R_f(U)\) iff \(c\in U\), and the complex nullary value is
\(\widehat c\).
The choices \(c=1\) and \(c=0\) give the whole Stone space and the
empty set, respectively.
\end{hbexample}

\begin{hblemma}{thm:p10-bao-ultrafilter-witness}{proved=annex:p10.bao.1}{Simultaneous ultrafilter witnesses}
Let \(f:B^n\to B\) be an operator, \(U\in X_B\), and
\(F_1,\ldots,F_n\) proper filters such that
\[
 f(a_1,\ldots,a_n)\in U
              \quad\text{whenever }a_i\in F_i\ (1\leq i\leq n).
\]
There are ultrafilters \(V_i\supseteq F_i\) such that
\(R_f(U,V_1,\ldots,V_n)\).
In particular, if \(f(a_1,\ldots,a_n)\in U\), the \(V_i\) can be
chosen to contain the respective \(a_i\).
For \(n=0\), the conclusion has the empty witness tuple and is
exactly the assumption \(f()\in U\).
\end{hblemma}

Normality is needed for the element-witness conclusion.
The unary constant-one operation on the two-element Boolean
algebra preserves binary joins and sends \(0\) into its ultrafilter.
It cannot have a witness ultrafilter containing that input \(0\).

\begin{hbtheorem}{thm:p10-jonsson-tarski}{proved=annex:p10.bao.2}{J\'onsson--Tarski representation}
The Stone map
\[
 \beta_B:B\longrightarrow(\mathbf B^\bullet)^+,
 \qquad b\longmapsto\widehat b,
\]
is an embedding of Boolean algebras with operators:
\[
 \widehat{f^B(b_1,\ldots,b_n)}
      =f^{(\mathbf B^\bullet)^+}(\widehat b_1,\ldots,\widehat b_n).
\]
This includes every finite arity and every nullary operation.
The image of the Boolean reduct is
\(\operatorname{Clop}(X_B)\).
For a trivial \(\mathbf B\), its ultrafilter frame is empty and
its complex algebra is trivial.
\end{hbtheorem}

The source states this theorem for the full finitary signature but
gives its operator-preservation proof for one unary diamond
\cite[preprint Theorem~5.15, p.~20]{Venema2005Preprint}.
The annex proves the simultaneous filter-extension lemma and all
finite-arity cases.

\subsection{Canonical extensions}
\label{sec:p10-bao-canonical-extension}

\begin{hbdefinition}{def:p10-bao-canonical-extension}{The canonical extension and its Boolean completion}
The canonical extension of \(\mathbf B\) is
\[
 \mathbf B^\sigma=(\mathbf B^\bullet)^+,
\]
with the embedding \(\beta_B\).
This is the source's canonical-extension construction
\cite[preprint Definition~5.16, p.~21]{Venema2005Preprint}.
Its underlying Boolean algebra is complete and atomic.

For a Boolean embedding \(e:B\to D\) into a complete Boolean
algebra, call a meet of elements of \(e[B]\) \emph{closed} and a
join of them \emph{open}, including empty meets and joins.
The embedding is \emph{dense} when every element of \(D\) is
both a join of closed elements and a meet of open elements.
It is \emph{compact} when
\[
 \bigwedge e[S]\leq\bigvee e[T]\quad(S,T\subseteq B)
\]
implies the corresponding inequality for some finite
\(S_0\subseteq S\), \(T_0\subseteq T\).
Equivalently, an inequality between the meet of a family of
closed elements and the join of a family of open elements has
a witness using finitely many members of those two families.
The annex proves the equivalence of these compactness formulations.
All infinite meets and joins here are taken in the completion \(D\).
These are order-theoretic completion conditions
\cite[preprint \S7.2, Definition~7.1]{Venema2005Preprint}.
\index{canonical extension}\index{Boolean completion}
\end{hbdefinition}

\begin{hbtheorem}{thm:p10-bao-canonical-completion}{proved=annex:p10.bao.3}{The dense and compact canonical completion}
The Boolean embedding
\(\beta_B:B\to\mathcal P(X_B)\) is dense and compact.
Any other dense and compact Boolean completion of \(B\) is
isomorphic to it by a unique Boolean isomorphism commuting with
the embeddings.
This uniqueness concerns the Boolean completion.
The operators of \(\mathbf B^\sigma\) are the explicitly defined
complex operators; each preserves arbitrary joins in every
positive-arity coordinate.
\end{hbtheorem}

The existence and uniqueness statement is the source's
Theorem~7.2, with the Boolean atomicity argument supplied in
the annex \cite[preprint pp.~38--39]{Venema2005Preprint}.
The embedding need not preserve an infinite join that already
exists in \(B\).
For example, let \(B\) be the Boolean algebra of finite and
cofinite subsets of \(\mathbb N\).
The join of all singleton elements is \(\mathbb N=1\) in \(B\).
The cofinite sets form an ultrafilter \(U_\infty\), since exactly
one of any finite/cofinite set and its complement is cofinite.
Then \(U_\infty\in\widehat1\) but
\(U_\infty\notin\bigcup_n\widehat{\{n\}}\).
Thus that union, the join of the singleton images in the
completion, is strictly below the image of their join in \(B\).

The embedding transfers a valuation in \(\mathbf B\) to one whose
values lie in \(\beta_B[B]\).
Consequently every algebraic equation valid in \(\mathbf B\)
holds for those transported valuations.
Validity of an additional equation for all subsets of \(X_B\)
requires a further argument; it is the issue addressed by
canonicity in the following modal-completeness development.

\subsection{Admissible sets and descriptive general frames}
\label{sec:p10-bao-general-frames}

\begin{hbdefinition}{def:p10-bao-general-frame}{A general frame for a finitary operator signature}
A general frame is
\(\mathcal G=(W,(R_f)_{f\in\Lambda},\mathcal A)\), where
\(\mathcal A\subseteq\mathcal P(W)\) is a Boolean subalgebra
closed under every complex operator \(f^{\mathfrak F^+}\).
Its elements are the admissible proposition sets.
Valuations in a general frame assign proposition variables to
members of \(\mathcal A\); recursive formula evaluation remains
in \(\mathcal A\).
Write \(\mathcal G^+\) for the Boolean algebra with operators
whose carrier is precisely \(\mathcal A\).
The underlying relational frame's full complex algebra has
carrier \(\mathcal P(W)\).
For the unary nonempty case this agrees with
Definition~\ref{def:p08-general-frame}
\cite[preprint Definition~5.24, p.~22]{Venema2005Preprint}.
\end{hbdefinition}

\begin{hbdefinition}{def:p10-bao-descriptive}{Descriptive general frames}
A general frame is \emph{differentiated} when its admissible sets
separate distinct points.
It is \emph{compact} when every family of admissible sets with
nonempty finite intersections has nonempty total intersection.
Here the empty finite intersection is \(W\); this convention
also permits the empty frame.
It is \emph{tight} when for every \(f\) and tuple,
\[
 R_f(w,w_1,\ldots,w_n)
 \quad\Longleftrightarrow\quad
 \forall A_i\in\mathcal A\,
 \left(\bigwedge_{i=1}^n w_i\in A_i
       \ \Longrightarrow\
       w\in f^{\mathcal G^+}(A_1,\ldots,A_n)\right).
\]
The nullary case has no \(A_i\) to quantify over.
A \emph{descriptive} general frame has all three properties.
Its topology consists of unions of admissible sets
\cite[preprint Definition~5.25, p.~22]{Venema2005Preprint}.
\end{hbdefinition}

\begin{hbtheorem}{thm:p10-bao-descriptive-representation}{proved=annex:p10.bao.4}{Descriptive representation and reconstruction}
For every \(\mathbf B\),
\[
 \mathbf B_*=(X_B,(R_f),\{\widehat b:b\in B\})
\]
is descriptive, and \(\beta_B:\mathbf B\to(\mathbf B_*)^+\)
is an isomorphism of Boolean algebras with operators.
Conversely, for every descriptive \(\mathcal G\), its induced
topology is a Stone topology, its admissible sets are exactly
the clopens, and
\[
 w\longmapsto\{A\in\mathcal A:w\in A\}
\]
is an isomorphism of general frames
\(\mathcal G\to(\mathcal G^+)_*\).
It is a homeomorphism, preserves and reflects every relation,
and identifies the admissible families.
\end{hbtheorem}

\begin{hbdefinition}{def:p10-bao-general-morphism}{Bounded morphisms respecting admissible sets}
A morphism \(g:\mathcal G\to\mathcal H\) of general frames is a
function on their carriers such that preimages of admissible
sets are admissible, and for each \(f\) it has both properties:
\[
\begin{aligned}
 R_f^{\mathcal G}(w,u_1,\ldots,u_n)
   &\Longrightarrow
     R_f^{\mathcal H}(g(w),g(u_1),\ldots,g(u_n)),\\
 R_f^{\mathcal H}(g(w),v_1,\ldots,v_n)
   &\Longrightarrow
     \exists u_1,\ldots,u_n\
       \left(R_f^{\mathcal G}(w,u_1,\ldots,u_n)
                 \land\bigwedge_i g(u_i)=v_i\right).
\end{aligned}
\]
For a nullary symbol these say
\(R_f^{\mathcal G}(w)\) iff \(R_f^{\mathcal H}(g(w))\).
The admissible-preimage condition makes the map continuous for
the induced topologies
\cite[preprint Definition~5.27, p.~23]{Venema2005Preprint}.
\end{hbdefinition}

\begin{hbtheorem}{thm:p10-bao-descriptive-duality}{proved=annex:p10.bao.5}{Duality with descriptive general frames}
Boolean algebras with the fixed \(\Lambda\)-operators and their
homomorphisms are contravariantly equivalent to descriptive
\(\Lambda\)-general frames with the preceding morphisms.
A homomorphism \(h:\mathbf A\to\mathbf B\) gives
\[
 h_*:\mathbf B_*\to\mathbf A_*,
 \qquad h_*(U)=h^{-1}[U],
\]
and a general-frame morphism gives the inverse-image
homomorphism on its admissible algebras.
The reconstruction maps of
Theorem~\ref{thm:p10-bao-descriptive-representation} are natural.
These statements include trivial algebras and empty descriptive
frames \cite[preprint Proposition~5.18 and Theorem~5.28,
pp.~21, 23]{Venema2005Preprint}.
\end{hbtheorem}
